\documentclass[11pt,a4paper]{article}

\usepackage[top=0.6in,bottom=0.6in,left=0.65in,right=0.65in]{geometry}
\usepackage[T1]{fontenc}
\usepackage{lmodern}
\usepackage{enumitem}
\usepackage{titlesec}
\usepackage{hyperref}
\usepackage{xcolor}
\usepackage{microtype}
\usepackage{tabularx}
\usepackage{needspace}
\usepackage{graphicx}
\usepackage{tikz}

% photo version for Swiss applications; set \photofalse for international roles
\newif\ifphoto
\phototrue

\definecolor{accent}{HTML}{1a1a1a}
\definecolor{subtle}{HTML}{333333}

\hypersetup{
  colorlinks=true,
  urlcolor=accent,
  linkcolor=accent,
  pdfauthor={Mariia Eremina},
  pdftitle={Mariia Eremina - Resume},
}

\pagestyle{empty}

\titleformat{\section}{\large\bfseries\color{accent}\uppercase}{}{0em}{}[\titlerule]
\titlespacing*{\section}{0pt}{14pt}{6pt}

\setlist[itemize]{leftmargin=1.4em,itemsep=2pt,parsep=0pt,topsep=3pt}

\newcommand{\role}[4]{%
  \needspace{5\baselineskip}%
  \noindent\textbf{#1}\hfill\textit{#2}\\
  \textbf{\textit{#3}}\hfill\textcolor{subtle}{#4}\par\vspace{2pt}%
}

\begin{document}

% ── HEADER ──────────────────────────────────────────────────
\newcommand{\contactlines}{%
  {\color{subtle}Basel, Switzerland\enspace\textperiodcentered{}\enspace+41\,76\,229\,4344\enspace\textperiodcentered{}\enspace%
  \href{mailto:mariia.erem@gmail.com}{mariia.erem@gmail.com}}\\[2pt]
  {\color{subtle}%
  \href{https://mariia.ch}{mariia.ch}\enspace\textperiodcentered{}\enspace%
  \href{https://linkedin.com/in/mariia-eremina}{linkedin.com/in/mariia-eremina}\enspace\textperiodcentered{}\enspace%
  \href{https://github.com/marerem}{github.com/marerem}}%
}
\ifphoto
\noindent
\begin{minipage}[c]{\dimexpr\textwidth-3cm\relax}
  {\LARGE\bfseries Mariia Eremina}\\[6pt]
  \contactlines
\end{minipage}\hfill
\begin{minipage}[c]{2.6cm}
  \raggedleft
  \begin{tikzpicture}
    \clip (0,0) circle (1.25cm);
    \node at (0.12,-0.2) {\includegraphics[width=3.4cm]{public/photos/portrait.jpg}};
  \end{tikzpicture}
\end{minipage}
\else
\begin{center}
  {\LARGE\bfseries Mariia Eremina}\\\vspace{4pt}
  \contactlines
\end{center}
\fi

\vspace{-2pt}

% ── SUMMARY ─────────────────────────────────────────────────
\section{Summary}

Full-stack ML engineer with nearly 4 years of industry experience taking ML from research to production: the models, the infrastructure that runs them, and the interfaces people use.
Built observability and orchestration platforms at MDPI, clinical ML used at hospitals across the Boston area at MIT, and production anomaly detection at Maxon; co-founded an EPFL-backed startup that reached $\sim$5,000 users.

% ── EXPERIENCE ──────────────────────────────────────────────
\section{Professional Experience}

\role{Independent AI Engineer}{May 2026 -- Present}{Personal projects}{Basel, Switzerland}
\begin{itemize}
  \item Deployed \href{https://www.gabyluong.com}{gabyluong.com} to production for a Senior UI/UX Designer.
  \item Self-hosted a Jev-style decision readout (Noul, Score, Choice) served with vLLM and PagedAttention on Kubernetes.
  \item Built \href{https://github.com/marerem/longmem}{longmem}, open-source cross-project memory for AI assistants (hybrid semantic + keyword search), released as an MCP server on Glama and published on PyPI.
\end{itemize}

\role{ML \& Infrastructure Engineer}{May 2025 -- Present}{MDPI, Open-Access Scientific Journal Publisher}{Basel, Switzerland}
\begin{itemize}
  \item Led the observability platform replacing Grafana--Prometheus (OpenTelemetry, ClickHouse, SigNoz), with alerting to Teams and Slack, \textbf{cutting incident-response time by 30\%}.
  \item Architected and built the unified API layer (FastAPI, PostgreSQL), with versioned Parquet backups on S3.
  \item Migrated ML and business-automation pipelines from Airflow to \textbf{self-hosted Prefect with zero downtime} and failure alerts.
  \item Moved the Python toolchain to Rust-based tools (Poetry to uv, pandas to Polars, Loguru to Logly) for faster dependency resolution, lower-memory ETL and structured logs.
  \item Managed teams across Switzerland and China and set architectural standards, including dbt for data transformations and lineage.
\end{itemize}

\role{Software Engineer}{Sep 2024 -- May 2025}{MIT, Massachusetts Institute of Technology}{Boston, USA (Hybrid)}
\begin{itemize}
  \item Took \textbf{4 of the lab's research prototypes to production} as one full-stack platform (React, Node.js), used by researchers and clinicians at \textbf{hospitals across the Boston area}.
  \item Served the PyTorch models through containerized APIs (Docker, Linux) and designed the interface from clinicians' user stories.
  \item Onboarded and mentored 2 new engineers.
\end{itemize}

\role{Co-Founder}{May 2024 -- Dec 2024}{CarGo Relay, C2C Parcel Delivery (EPFL Startup)}{Lausanne, Switzerland}
\begin{itemize}
  \item 1 of 3 co-founders of a C2C parcel-delivery startup backed by EPFL Startup Support, matching parcels to the routes people travel most from Monday to Friday.
  \item Grew to \textbf{$\sim$5,000 registered users} across universities and schools in Lausanne and \textbf{$\sim$CHF 50,000 in revenue} by the end of 2024 from a 10\% commission on deliveries.
  \item Led system design and built the product end to end (React Native, FastAPI, MariaDB, Google Maps API), including mode-aware route matching: close to the route for metro riders, wider detours for bikes, and large parcels to cars.
\end{itemize}

\role{Research Engineer}{Mar 2024 -- Sep 2024}{MIT \& Harvard Medical School, Edelman Lab}{Boston, USA}
\begin{itemize}
  \item Invented a landmark-based co-registration algorithm that aligns video frames by spline interpolation between landmarks, replacing a slow, GPU-bound learned model that overfit on scarce clinical data.
  \item Made co-registration \textbf{40$\times$ faster} on CPU, running inside the clinical workflow on the imaging machine itself.
  \item Built a frame-cleaning pipeline (CLIPSeg and SAM segment the medical equipment, Stable Diffusion inpainting removes it) and made it cross-platform (macOS, Windows, Linux) for medical hardware.
\end{itemize}

\role{Software \& Machine Learning Engineer}{Jun 2023 -- Mar 2024}{Maxon Group, Precision Motor Systems}{Lucerne, Switzerland}
\begin{itemize}
  \item Partnered with 2 AI experts from CSEM and 3 engineers to detect defective precision motors from their current, voltage and temperature signals, where defects were rare and masked by factory noise.
  \item Benchmarked a classical pipeline (signal feature extraction, dimensionality reduction, gradient boosting) against deep models (autoencoder, Siamese network), optimizing for recall since a missed defect costs more than a false alarm.
  \item Shipped the winning autoencoder (\textbf{92\% recall}) as an inspection application (FastAPI, PostgreSQL, Docker, JWT auth) integrated into the factory production line, with \textbf{90\%+ test coverage} in CI.
\end{itemize}

\role{Machine Learning Engineer}{Dec 2022 -- May 2023}{AGORA Cancer Research Center}{Lausanne, Switzerland}
\begin{itemize}
  \item Built a detection-and-tracking model (YOLOv5, DeepSORT) that follows each cancer organoid as it grows, supporting clinical research at CHUV.
  \item Classified organoids into 5 readiness stages for drug testing, oversampling rare stages and weighting the loss by class frequency, since debris and non-viable organoids dominated the data.
  \item Replaced fully manual, frame-by-frame annotation with a human-in-the-loop pipeline that \textbf{automated 89\% of predictions}, sending only uncertain ones for human approval.
  \item Architected and maintained the lab's 150\,TB imaging data pipeline.
\end{itemize}

% ── EDUCATION ───────────────────────────────────────────────
\section{Education}

\role{Master's Thesis in Computer Science, Edelman Lab}{2023 -- 2024}{MIT, Massachusetts Institute of Technology}{Boston, USA}
\vspace{-2pt}
\begin{itemize}
  \item ``Landmark-Based Co-Registration of Coronary CT and Intravascular Images''
\end{itemize}

\role{MSc Computer Science \& Life Science, Minor in Communication Systems}{2021 -- 2024}{EPFL, Swiss Federal Institute of Technology}{Lausanne, Switzerland}
\vspace{-2pt}
\begin{itemize}
  \item Organise annual hackathons sponsored by AWS, UBS, Logitech, SBB, Swisscom, Swissquote, and Huawei (ongoing since 2022).
\end{itemize}

\role{BSc Biotechnology}{2017 -- 2021}{Lomonosov Moscow State University}{Moscow, Russia}

% ── SKILLS ──────────────────────────────────────────────────
\section{Technical Skills}

\begin{tabularx}{\textwidth}{@{}l X@{}}
\textbf{Infrastructure \& DevOps} & Docker, Kubernetes, AWS (S3), GCP, Azure, Linux, CI/CD (GitHub Actions, GitLab CI), Ansible, Prefect, dbt, uv, Git \\[3pt]
\textbf{Observability} & OpenTelemetry, ClickHouse, SigNoz, Prometheus, Grafana, Logly \\[3pt]
\textbf{ML \& Data} & PyTorch, TensorFlow, Keras, OpenCV, scikit-learn, vLLM, MCP, Hydra, Polars, pandas, NumPy, Optuna, W\&B \\[3pt]
\textbf{Programming} & Python, TypeScript, JavaScript, HTML, SQL, R, MATLAB, \LaTeX \\[3pt]
\textbf{Web \& Backend} & React, React Native, Node.js, FastAPI, Flask, PostgreSQL, MariaDB, MongoDB, SQLite, Redis, Firebase \\[3pt]
\textbf{Spoken Languages} & English (C1), Russian (native), German (A1) \\
\end{tabularx}

\end{document}
